\documentclass[11pt,a4paper]{article}
\usepackage[T1]{fontenc}
\usepackage[utf8]{inputenc}
\usepackage{amsmath,amsthm,mathtools}
\usepackage{newtxtext,newtxmath}
\usepackage[margin=26mm,headheight=15pt]{geometry}
\usepackage{microtype}
\usepackage{bm}
\usepackage{booktabs,longtable,array,tabularx}
\usepackage{xcolor,graphicx}
\usepackage{enumitem}
\usepackage{listings}
\usepackage{needspace}
\usepackage{xurl}
\usepackage{fancyhdr}
\usepackage[colorlinks=true,linkcolor=blue!45!black,citecolor=blue!45!black,urlcolor=blue!45!black]{hyperref}
\usepackage[nameinlink,noabbrev]{cleveref}
\hypersetup{pdftitle={Positive-Real Asymptotic Inversion with Powers and Logarithms},pdfauthor={Prepared with ChatGPT for Vladimir Reshetnikov},pdfsubject={Constructive power-log inverse expansions and Wolfram Language implementation}}
\definecolor{ink}{RGB}{28,43,59}
\definecolor{soft}{RGB}{244,247,250}
\definecolor{ruleblue}{RGB}{49,87,119}
\pagestyle{fancy}
\fancyhf{}
\fancyhead[L]{\small\sffamily PowerLogInverse\enspace /\enspace Theory and implementation}
\fancyhead[R]{\small\sffamily 1.0.0}
\fancyfoot[C]{\small\thepage}
\renewcommand{\headrulewidth}{0.3pt}
\setlength{\parindent}{1.1em}
\setlength{\parskip}{3pt}
\setlength{\emergencystretch}{2em}
\allowdisplaybreaks[2]
\newtheorem{theorem}{Theorem}[section]
\newtheorem{proposition}[theorem]{Proposition}
\newtheorem{lemma}[theorem]{Lemma}
\newtheorem{corollary}[theorem]{Corollary}
\theoremstyle{definition}
\newtheorem{definition}[theorem]{Definition}
\newtheorem{example}[theorem]{Example}
\theoremstyle{remark}
\newtheorem{remark}[theorem]{Remark}
\newcommand{\OO}{\mathcal O}
\newcommand{\NN}{\mathbb N_0}
\newcommand{\RR}{\mathbb R}
\newcommand{\CC}{\mathbb C}
\newcommand{\dd}{\mathrm d}
\newcommand{\ee}{\mathrm e}
\newcommand{\val}{\operatorname{val}}
\newcommand{\coeff}{\operatorname{coeff}}
\newcommand{\Li}{\mathscr L}
\newcommand{\T}{\mathcal T}
\newcommand{\fall}[2]{#1^{\underline{#2}}}
\newcommand{\code}[1]{\texttt{#1}}
\newcommand{\pkg}{\textsf{PowerLogInverse}}
\lstdefinelanguage{WL}{morekeywords={Module,Block,With,Function,Table,Total,Expand,D,Get,Clear,If,Return,Failure,Sum,FullSimplify,Assumptions,Log,Sqrt,PowerLogInverse,PowerLogInverseData,PowerLogModel,PowerLogCoefficient,TestReport,VerificationTest,Options,OptionsPattern,OptionValue,True,False,Normal,Series,InverseSeries,Collect,N,Abs,Exp,Print,FileNameJoin,DirectoryName},sensitive=true,morecomment=[s]{(*}{*)},morestring=[b]"}
\lstset{language=WL,basicstyle=\ttfamily\small,keywordstyle=\color{ruleblue},commentstyle=\itshape\color{black!55},stringstyle=\color{black!75},backgroundcolor=\color{soft},frame=single,rulecolor=\color{black!15},framerule=0.3pt,breaklines=true,columns=fullflexible,keepspaces=true,showstringspaces=false,tabsize=2,xleftmargin=3pt,xrightmargin=3pt,aboveskip=9pt,belowskip=9pt}
\setcounter{tocdepth}{2}
\makeatletter
\newcommand{\tableinput}[1]{\@@input #1 }
\makeatother

\begin{document}
\begin{titlepage}
\thispagestyle{empty}
\color{ink}
{\sffamily\small MATHEMATICAL SOFTWARE\enspace /\enspace RESEARCH REPORT}\par
\vspace{12mm}
{\Huge\bfseries Positive-Real\par Asymptotic Inversion\par}
\vspace{4mm}
{\LARGE with Powers and Logarithms\par}
\vspace{7mm}
{\large A constructive coefficient theory\par and a Wolfram Language package\par}
\vspace{8mm}
\noindent\color{ruleblue}\rule{\textwidth}{1pt}\color{ink}
\vspace{5mm}
{\large Prepared for Vladimir Reshetnikov\par}
{\normalsize Prepared with ChatGPT\par}
\vspace{3mm}
{\sffamily Version 1.0.0\enspace\textbullet\enspace September 7, 2026\par}
\vspace{10mm}
\begin{abstract}
\noindent We construct the positive real inverse near zero of
\[
 f(x)=a x^p\left(1+\sum_{j=1}^{m}x^{\delta_j}P_j(\log x)\right),
 \qquad a,p,\delta_j>0,
\]
where the $P_j$ are real polynomials and the gaps $\delta_j$ may be irrational.
A homotopy form of Lagrange--B\"urmann inversion produces every coefficient by a
finite product of first-order differential operators on polynomials. We prove
local existence, branch selection, convergence of the polynomial-block
expansion, formal uniqueness, and logarithm-aware remainder estimates for both
exponent and perturbation-depth truncation. The original examples
$x+x^2(1+\log x)$ and $x+x^{\sqrt2}$ are worked out to arbitrary order, with
explicit coefficients and convergence information. The accompanying package
normalizes admissible input, preserves exact irrational exponents, combines
coincident weights, and returns expressions together with error metadata.
\end{abstract}
\vfill
\noindent\fcolorbox{ruleblue}{soft}{\parbox{0.94\textwidth}{\small
\textbf{Verification status.} The mathematical specification passed
\input{verification-count.tex}\ independent executed checks, including formal
substitution and 100-digit numerical comparisons. The Wolfram source and 35
native regression tests are supplied, but native kernel execution was
unavailable in this session. No native test success is claimed.}}
\end{titlepage}

\tableofcontents
\newpage

\section{The question, the answer, and the correct real branch}
\label{sec:question}

The archived question asks for an asymptotic inverse on the positive real axis,
not for an arbitrary local complex inverse. Its principal example is
\begin{equation}\label{eq:original}
 f(x)=x+x^2(1+\log x),\qquad x>0,
\end{equation}
and its second example is $x+x^{\sqrt2}$. The reported computer-algebra attempts
sometimes remain unevaluated and sometimes select a non-real inverse branch
\cite{question}. The mathematical issue is not resolved merely by appending
$y>0$ to an otherwise unspecified inverse: one must choose the inverse germ
whose argument tends to zero along the positive real axis.

For \eqref{eq:original}, the answer begins
\begin{align}
 g(y)={}&y-y^2(1+\log y)
       +y^3\bigl(2\log^2y+5\log y+3\bigr)\notag\\
 &-y^4\left(5\log^3y+\frac{41}{2}\log^2y
                     +27\log y+\frac{23}{2}\right)
       +\OO\bigl(y^5|\log y|^4\bigr).
 \label{eq:log-answer}
\end{align}
All limits in this report are $y\to0^+$ unless another regime is explicitly
specified. The exact all-orders coefficient formula is given in
\cref{sec:logexample}; it is not inferred by numerical fitting.

For the irrational-power example the answer is
\begin{align}
 (x+x^{\sqrt2})^{-1}(y)
 ={}&y-y^{\sqrt2}+\sqrt2\,y^{2\sqrt2-1}
 -\frac{6-\sqrt2}{2}\,y^{3\sqrt2-2}\notag\\
 &+\left(\frac{17\sqrt2}{3}-4\right)y^{4\sqrt2-3}
 +\OO\bigl(y^{5\sqrt2-4}\bigr).
 \label{eq:irrational-answer}
\end{align}
Here $f^{-1}$ means functional inversion, not a reciprocal. The first four
terms reproduce the expected result in the supplied archive, and the next
term is generated by the same rule as every subsequent term.

\subsection{Why the positive inverse exists globally in the two examples}

For \eqref{eq:original},
\[
 f'(x)=1+x(3+2\log x).
\]
The derivative of $x(3+2\log x)$ is $5+2\log x$, so its global minimum occurs at
$x=\ee^{-5/2}$. Consequently
\begin{equation}\label{eq:global-slope}
 f'(x)\ge m_0:=1-2\ee^{-5/2}>0\qquad(x>0).
\end{equation}
Also $f(x)\to0$ as $x\to0^+$ and $f(x)\to+\infty$ as $x\to+\infty$.
Thus $f$ is a bijection of $(0,\infty)$ onto itself. Its inverse satisfies
$g(y)\sim y$. The sign of the first correction is useful as a check: for small
$y$, $1+\log y<0$, and therefore $g(y)>y$.

For $f_\alpha(x)=x+x^\alpha$, with any fixed $\alpha>1$,
\[
 f_\alpha'(x)=1+\alpha x^{\alpha-1}\ge1,
\]
so the same global positive bijection conclusion holds. In this case
$g_\alpha(y)<y$. Both inequalities agree with the expansions.

The general model studied below need only be monotone sufficiently near zero.
Higher-power terms with negative coefficients can destroy global monotonicity
without affecting the positive local inverse selected by the package.

\subsection{Why an ordinary Taylor branch can be the wrong object}

Away from the logarithmic cut, the nonzero complex solutions of $f(z)=0$ in
\eqref{eq:original} obey
\[
 z(1+\log z)=-1.
\]
Writing $w=1+\log z$ gives $w\ee^w=-\ee$. Thus such roots can be described,
with compatible logarithm branches, by
\[
 z_0=\exp\bigl(W_k(-\ee)-1\bigr).
\]
At one of these roots $f'(z_0)=z_0-1\ne0$, so a perfectly legitimate analytic
inverse about $y=0$ has the expansion
\[
 z_0+\frac{y}{z_0-1}
 -\frac{3+2W_k(-\ee)}{2(z_0-1)^3}y^2+\cdots.
\]
It is not the positive inverse tending to zero. This calculation explains the
shape of the non-real expansion reported in the archive; it does not diagnose
a current-version software bug.

The desired inverse is not analytic at the endpoint. Indeed,
$f''(x)=5+2\log x\to-\infty$, and
\[
 g''(y)=-\frac{f''(g(y))}{f'(g(y))^3}\longrightarrow+\infty.
\]
It cannot have an ordinary Taylor expansion at zero. The failure of ordinary
Taylor analyticity does not preclude a precise, convergent expansion in
power--logarithm blocks.

\subsection{What is contributed here}

The coefficient identity is a specialization of classical
Lagrange--B\"urmann reversion, not a claim of a new general inversion theorem.
Classical inverse-function and residue formulations are recorded in
\cite{dlmf}. The constructive contribution is to normalize a broad class of
positive real germs, express its coefficients entirely through finite
polynomial operations, justify exact irrational-weight truncation, and
implement the resulting specification with explicit failure modes.

The supplied repository is useful methodological context. Its directory
README and inspected portions of the companion monograph emphasize choosing
a dominant inverse core, organizing perturbations by weight, and transporting
a forward residual through a slope estimate \cite{repo,companion}. Here these
ideas are specialized to a class for which all analytic assertions can be
proved directly. No uninspected theorem in a larger transseries monograph is
used as a premise.

\section{Power--logarithm scales and admissible models}
\label{sec:scales}

\begin{definition}[Admissible finite model]\label{def:model}
An admissible model is
\begin{equation}\label{eq:model}
 f(x)=a x^p\bigl(1+R(x)\bigr),\qquad
 R(x)=\sum_{j=1}^{m}x^{\delta_j}P_j(\log x),
\end{equation}
where $a>0$, $p>0$, each $\delta_j>0$, and each $P_j$ is a nonzero real
polynomial. The number $m$ is finite; $m=0$ is allowed and gives a pure
monomial. All parameters are fixed. Powers and logarithms are real on
$x>0$, and their complex continuations are taken on small disks avoiding the
nonpositive real axis when needed in proofs.
\end{definition}

The input class includes polynomial perturbations, rational Puiseux powers,
exact irrational powers, arbitrary finite polynomial degrees in $\log x$,
and mixtures of these features. Coefficients of the $P_j$ may have either
sign. Repeated gaps are permitted mathematically and may be combined by adding
their polynomials.

Put
\begin{equation}\label{eq:normalization}
 t=\left(\frac ya\right)^{1/p},\qquad L=\log t,
 \qquad T=1+|L|.
\end{equation}
Then $t>0$ and $t\to0^+$ exactly when $y\to0^+$. The inverse sought is
$g(y)=t(1+o(1))$.

\subsection{The hierarchy of powers and logarithms}

For every fixed $\varepsilon>0$ and real $M$,
\begin{equation}\label{eq:power-beats-log}
 t^\varepsilon(1+|\log t|)^M\longrightarrow0.
\end{equation}
One proof is to write $t=\ee^{-u}$: an exponential decay beats any fixed power
of $u$. It follows that any finite logarithmic polynomial at a larger power of
$t$ is smaller than any nonzero logarithmic polynomial at a smaller power.
Within a fixed power, the term of highest logarithmic degree dominates in
absolute value.

This is why exponent blocks are natural:
\[
 t^\beta Q(L),\qquad Q\in\RR[L].
\]
A block is kept as a whole. There is no need to rank its finitely many
logarithmic monomials separately in the implementation.

A crucial error distinction is
\[
 t^\beta\log^2t\notin\OO(t^\beta),
 \qquad
 t^\beta\log^2t\in\OO(t^{\beta-\varepsilon})
 \quad\text{for every }\varepsilon>0.
\]
An order term that suppresses the logarithms can therefore assert a false
error bound even when the displayed coefficients are right.

\subsection{The positive additive support}

For $k=(k_1,\ldots,k_m)\in\NN^m$, write
\[
 |k|=\sum_j k_j,\qquad k!=\prod_j k_j!,\qquad
 \Delta_k=\sum_j k_j\delta_j.
\]
The generated weight set is
\[
 S=\{\Delta_k:k\in\NN^m\}.
\]
If $d=\min_j\delta_j>0$, then $\Delta_k\ge d|k|$. Thus there are only finitely
many multiindices with weight below any prescribed finite cutoff. In
particular, $S$ is locally finite, even when its generators are irrational.

It is important not to confuse this positive semigroup with the additive
group generated by the same numbers using negative integer coefficients. That
group can be dense, while the positive semigroup relevant here has only
finitely many points below any bound.

Distinct multiindices can yield the same weight. For example, gaps $1$ and $2$
give $\Delta_{(2,0)}=\Delta_{(0,1)}=2$. The coefficient at weight $2$ is the sum
of both contributions. Declaring the two terms to be different asymptotic
scales would be incorrect. Conversely, replacing irrational gaps by decimal
approximations can create spurious equalities. The package does neither.

\subsection{Existence of the general positive inverse}

\begin{proposition}\label{prop:localinverse}
For an admissible model, $R(x)=o(1)$ and $xR'(x)=o(1)$. Consequently
\[
 f'(x)=apx^{p-1}(1+o(1)),
\]
and $f$ has a unique positive inverse for all sufficiently small $y>0$.
Moreover $g(y)\sim t$ with $t$ defined by \eqref{eq:normalization}.
\end{proposition}

\begin{proof}
Every $x^{\delta_j}P_j(\log x)$ tends to zero by
\eqref{eq:power-beats-log}. Also
\[
 x\frac{\dd}{\dd x}\bigl(x^{\delta_j}P_j(\log x)\bigr)
 =x^{\delta_j}\bigl(\delta_jP_j(\log x)+P_j'(\log x)\bigr)=o(1).
\]
Differentiate \eqref{eq:model}. The factor multiplying $ax^{p-1}$ is
$p(1+R)+xR'$, which tends to $p>0$. Thus $f'>0$ and $f>0$ on a sufficiently
small interval $(0,x_0)$. The endpoint limit is zero, so $f$ maps this interval
bijectively onto $(0,f(x_0))$. Finally $y=ag(y)^p(1+o(1))$, giving $g(y)/t\to1$.
\end{proof}

\section{The coefficient theorem}
\label{sec:theorem}

Let $D_L=\dd/\dd L$. Fix $q>0$; the desired inverse itself corresponds to
$q=1$. Allowing $q$ also computes powers of the inverse without first making an
insufficiently long truncation of $g$ and then raising it to a power.

For $k\ne0$, define
\begin{equation}\label{eq:Qk}
 \boxed{\displaystyle
 Q_k^{(q)}(L)=
 \frac{(-1)^{|k|}q}{p^{|k|}k!}
 \left[\prod_{r=1}^{|k|-1}
       \bigl(q+\Delta_k+pr+D_L\bigr)\right]
 \prod_{j=1}^{m}P_j(L)^{k_j}.}
\end{equation}
An empty operator product is the identity. Set $Q_0^{(q)}=1$.
All operators in the product commute because their scalar parts are constants
with respect to $L$.

\begin{theorem}[Constructive positive-real inversion]\label{thm:main}
For every admissible finite model and every fixed $q>0$,
\begin{equation}\label{eq:inverse-main}
 g(y)^q=\sum_{k\in\NN^m}t^{q+\Delta_k}Q_k^{(q)}(\log t)
\end{equation}
for sufficiently small positive $t$. The series is absolutely convergent when
each polynomial $Q_k^{(q)}$ is regarded as one block. All multiindices at the
same weight may therefore be collected. In addition,
\begin{equation}\label{eq:degreebound}
 \deg Q_k^{(q)}\le\sum_j k_j\deg P_j.
\end{equation}
Every finite exponent cutoff can be computed by finitely many exact
polynomial operations and comparisons of real weights. The resulting
expansion is the positive inverse branch specified by $g(y)\sim t$.
\end{theorem}

The theorem has no rationality requirement on $p$ or the $\delta_j$.
For software, exact comparison may require assumptions or may remain
undecidable to a particular simplifier; that computational issue does not
alter the mathematical statement.

\subsection{The first two perturbation depths}

Write $P_j=P_j(L)$ and $P_j'=D_LP_j$. Formula \eqref{eq:Qk} yields
\begin{align}
 g(y)^q={}&t^q-\frac qp\sum_j t^{q+\delta_j}P_j\notag\\
 &+\frac{q}{2p^2}\sum_j t^{q+2\delta_j}
   \bigl((q+2\delta_j+p)P_j^2+2P_jP_j'\bigr)\notag\\
 &+\frac q{p^2}\sum_{i<j}t^{q+\delta_i+\delta_j}
   \bigl((q+\delta_i+\delta_j+p)P_iP_j+P_i'P_j+P_iP_j'\bigr)
   +\text{higher depths}.
 \label{eq:first-two}
\end{align}
The cross terms are essential. Independently inverting each perturbation and
adding the answers would miss them.

\subsection{A polynomial-vector implementation}

If $Z(L)=\sum_{s=0}^{K}z_sL^s$, then
\begin{equation}\label{eq:vector-operator}
 (A+D_L)Z=\sum_{s=0}^{K}\bigl(Az_s+(s+1)z_{s+1}\bigr)L^s,
 \qquad z_{K+1}=0.
\end{equation}
Thus every operator in \eqref{eq:Qk} acts by a short upper-bidiagonal
transformation on the coefficient vector. No unknown coefficient system has
to be solved. The supplied Wolfram implementation applies the equivalent
polynomial operation \code{Expand[A z + D[z,ell]]}; a specialized sparse-vector
backend is a possible optimization, not a prerequisite for the formula.

\section{Proof by homotopy reversion}
\label{sec:proof}

The essential device is to apply analytic reversion at a \emph{positive}
basepoint and expand in an auxiliary parameter. We never pretend that a
logarithmic function is analytic at the endpoint zero.

\subsection{The homotopy Lagrange--B\"urmann identity}

\begin{lemma}\label{lem:homotopy}
Let $H$ and $A$ be analytic on a disk about $v\ne0$. For sufficiently small
$\lambda$, let $s_\lambda$ be the solution near $v$ of
\[
 s_\lambda+\lambda H(s_\lambda)=v.
\]
Then
\begin{equation}\label{eq:homotopy}
 A(s_\lambda)=A(v)+\sum_{n\ge1}\frac{(-\lambda)^n}{n!}
 \left(\frac{\dd}{\dd v}\right)^{n-1}
 \bigl(A'(v)H(v)^n\bigr).
\end{equation}
\end{lemma}

\begin{proof}
Choose a small circle $C$ about $v$ on which
$|\lambda H(s)|<|s-v|$. Rouch\'e's theorem gives exactly one zero inside,
counted with multiplicity. It is simple, and the argument-principle residue
formula gives
\[
 A(s_\lambda)=\frac1{2\pi i}\int_C
 A(s)\frac{1+\lambda H'(s)}{s-v+\lambda H(s)}\,\dd s.
\]
The quotient has the uniformly convergent expansion on $C$
\[
 \frac{1+\lambda H'}{s-v+\lambda H}
 =\frac1{s-v}+\sum_{n\ge1}\frac{(-1)^{n+1}\lambda^n}{n}
       \frac{\dd}{\dd s}\left(\frac{H(s)^n}{(s-v)^n}\right).
\]
Integrate each coefficient by parts around the closed contour. The coefficient
of $\lambda^n$ becomes
\[
 \frac{(-1)^n}{n}\frac1{2\pi i}\int_C
       \frac{A'(s)H(s)^n}{(s-v)^n}\,\dd s.
\]
Cauchy's differentiation formula supplies the factor $1/(n-1)!$, proving
\eqref{eq:homotopy}. These are the usual analytic ingredients of Lagrange
reversion \cite{dlmf}; here the small parameter is $\lambda$, not $v$.
\end{proof}

For the near-identity case $f(x)=x+h(x)$, choose $A(s)=s$ and $v=y$. The
immediately useful formula is
\begin{equation}\label{eq:nearidentity}
 g(y)=y+\sum_{n\ge1}\frac{(-1)^n}{n!}
       \frac{\dd^{n-1}}{\dd y^{n-1}}h(y)^n,
\end{equation}
provided the homotopy expansion reaches $\lambda=1$. This proviso is proved
below for all admissible finite models.

\subsection{Removing the leading power}

Set
\[
 s=x^p,\qquad v=y/a=t^p,
\]
and define
\[
 H_j(s)=s^{1+\delta_j/p}P_j\!\left(\frac{\log s}{p}\right),
 \qquad H(s)=\sum_jH_j(s).
\]
The inverse equation becomes $s+H(s)=v$. To obtain $g(y)^q$ use
$A(s)=s^{q/p}$ in \eqref{eq:homotopy}.

Let $n=|k|$, $\Delta=\Delta_k$, and
$Z_k(L)=\prod_jP_j(L)^{k_j}$. The multinomial contribution indexed by $k$ in
$A'(s)H(s)^n$ is
\[
 \frac qp\frac{n!}{k!}\,
 s^{n-1+(q+\Delta)/p}Z_k\!\left(\frac{\log s}{p}\right).
\]
The elementary differentiation identity
\begin{equation}\label{eq:derivative-identity}
 \frac{\dd^r}{\dd s^r}\left[s^b Z\!\left(\frac{\log s}{p}\right)\right]
 =s^{b-r}\prod_{j=0}^{r-1}\left(b-j+\frac1pD_L\right)Z(L)
\end{equation}
follows by induction. Apply it with $r=n-1$ and
$b=n-1+(q+\Delta)/p$. Factoring $1/p$ from every operator gives exactly
\eqref{eq:Qk}, after $s=v=t^p$. This proves the coefficient identity formally
and locally in the homotopy parameter. The degree bound follows because
$D_L$ only lowers polynomial degree.

\subsection{Why the homotopy reaches one and the block series converges}

Fix $0<r<1$ and consider $|s-v|\le rv$, where $v=t^p>0$. A single analytic
logarithm exists on this disk. Since $s/v$ remains in the compact disk
$|s/v-1|\le r$,
\[
 \frac{\log s}{p}=L+\frac1p\log(s/v)
\]
differs from $L$ by a bounded quantity. For fixed model parameters there are
constants $C_j>0$, independent of sufficiently small $t$, such that
\begin{equation}\label{eq:Hbound}
 |H_j(s)|\le C_jv\,\varepsilon_j(t),\qquad
 \varepsilon_j(t)=t^{\delta_j}T^{K_j},\quad K_j=\deg P_j.
\end{equation}
Every $\varepsilon_j(t)$ tends to zero.

Introduce independent parameters $\lambda_j$ and solve
\[
 s+\sum_j\lambda_jH_j(s)=v.
\]
For $m>0$, choose the polydisk radii
\[
 \mathcal R_j(t)=\frac{r}{2mC_j\varepsilon_j(t)}.
\]
On $|s-v|=rv$ and $|\lambda_j|\le\mathcal R_j$ the sum of the perturbation
bounds is at most $rv/2$. Rouch\'e's theorem again gives exactly one simple root
inside. Local analytic dependence on the parameters patches to a single
analytic root throughout the polydisk because the root is unique.
Furthermore $|s/v-1|<r$, so the observable $s^{q/p}$ is analytic there and
bounded in absolute value by $C t^q$ for a fixed $C$.

The multivariate Cauchy coefficient estimate now gives
\begin{equation}\label{eq:multi-cauchy}
 \left|t^{q+\Delta_k}Q_k^{(q)}(L)\right|
 \le C t^q\prod_j\mathcal R_j(t)^{-k_j}.
\end{equation}
All radii tend to infinity. For sufficiently small $t$, the geometric sum of
the right-hand side over $k\in\NN^m$ is finite. Hence the Taylor series in the
$\lambda_j$ converges absolutely at $(1,\ldots,1)$ and gives
\eqref{eq:inverse-main}. Its coefficients agree with the multinomial
calculation above, for instance by first substituting
$\lambda_j=\lambda\mu_j$ and comparing powers of $\lambda$ and the $\mu_j$.

At $(1,\ldots,1)$ the root is real by complex-conjugation symmetry and unique;
it is positive because it lies in the disk about $v>0$ of radius less than $v$.
It is the local inverse of \cref{prop:localinverse}. This completes the proof
of \cref{thm:main}.

\begin{remark}
The convergence assertion groups each polynomial $Q_k^{(q)}(L)$ as one block.
It permits reordering and combining equal weights, which is all the algorithm
needs. It is not an assertion that an arbitrary ungrouping or resummation
convention for more general transseries preserves convergence.
\end{remark}

\section{Formal uniqueness and finite determination}
\label{sec:formal}

The analytic proof explains why the result is a genuine inverse. The formal
algebra explains why each requested finite part can be computed without
examining infinitely many terms.

Consider formal sums
\[
 \sum_{\Delta\in S}t^\Delta A_\Delta(L),\qquad A_\Delta\in\RR[L],
\]
with $S$ as in \cref{sec:scales}. Addition and multiplication are well defined
because each bounded weight range contains only finitely many indices.
The weight valuation is the smallest exponent with nonzero polynomial
coefficient. The Euler derivation acts by
\[
 t\frac{\dd}{\dd t}\bigl(t^\Delta A(L)\bigr)
   =t^\Delta(\Delta A+A'),
\]
so it preserves weights.

If $U$ has positive valuation, the binomial series $(1+U)^b$ and the formal
series $\log(1+U)$ are well defined: only finitely many powers of $U$ contribute
below any finite weight. Therefore substitution into a power--log term takes
the unambiguous form
\begin{equation}\label{eq:formalcomposition}
 (t(1+U))^\delta P(\log(t(1+U)))
 =t^\delta(1+U)^\delta P\bigl(L+\log(1+U)\bigr).
\end{equation}
The logarithm addition here is justified analytically by the chosen positive
branch and formally by the normalization of $1+U$; it is not an unrestricted
computer-algebra power expansion.

\begin{proposition}[Formal existence and uniqueness]\label{prop:formal}
There is a unique formal series $U$ with zero constant term such that
$g=t(1+U)$ solves the normalized inverse equation. Each coefficient below a
fixed weight depends on only finitely many input coefficients and finite
polynomial operations.
\end{proposition}

\begin{proof}
For $m=0$, the normalized solution is $U=0$. For $m>0$, the equation for $U$ is
\[
 (1+U)^p\left(1+\sum_jt^{\delta_j}(1+U)^{\delta_j}
   P_j\bigl(L+\log(1+U)\bigr)\right)=1.
\]
Equivalently,
\begin{equation}\label{eq:formalpicard}
 U=\left(1+\sum_jt^{\delta_j}(1+U)^{\delta_j}
   P_j\bigl(L+\log(1+U)\bigr)\right)^{-1/p}-1.
\end{equation}
Start with $U_0=0$ and repeatedly apply the right-hand side. A difference of
weight $w$ between two inputs produces a difference of weight at least $w+d$
between their outputs, because every perturbation carries a factor of weight
at least $d=\min\delta_j>0$. Thus the iteration stabilizes below any fixed
cutoff after finitely many steps and defines a solution.

For uniqueness, suppose two solutions first differ at weight $w$. In the
original normalized equation their difference at weight $w$ is $p$ times that
nonzero difference; all perturbation contributions have higher weight. Since
$p\ne0$, both equations cannot vanish. The finiteness assertion follows from
local finiteness of the support.
\end{proof}

The independent verification program uses the marker-variable version of
\eqref{eq:formalpicard}, not the differential coefficient formula, to check the
coefficients. This makes the verification substantially stronger than printing
the same formula in two programming languages.

\section{Remainders and two different truncation policies}
\label{sec:remainders}

The formulas are most useful when accompanied by a precise interpretation of
``order.'' The package supports two policies and reports which one was used.

\subsection{Truncation by total perturbation depth}

For $N\in\NN$, define
\[
 G_{q,N}^{\mathrm{depth}}(t)
 =\sum_{|k|\le N}t^{q+\Delta_k}Q_k^{(q)}(L).
\]
For $m>0$ put $d=\min_j\delta_j$ and $K=\max_j\deg P_j$.

\begin{theorem}[Depth remainder]\label{thm:depth}
For each fixed $N$,
\begin{equation}\label{eq:depthbound}
 g(y)^q-G_{q,N}^{\mathrm{depth}}(t)
 =\OO\left(t^{q+(N+1)d}T^{(N+1)K}\right).
\end{equation}
For $m=0$ the answer $g(y)^q=t^q$ is exact.
\end{theorem}

\begin{proof}
Apply \cref{lem:homotopy} with a single parameter multiplying $H=\sum H_j$.
On the disk used in \eqref{eq:Hbound},
\[
 |H(s)|\le C v\varepsilon(t),\qquad
 \varepsilon(t)=t^dT^K\longrightarrow0.
\]
Rouch\'e's theorem gives a homotopy convergence radius at least
$c/\varepsilon(t)$ with a fixed $c>0$. Cauchy's estimate bounds the coefficient
of depth $n$ by $C't^q(C''\varepsilon(t))^n$. Summing its geometric tail from
$n=N+1$ proves \eqref{eq:depthbound}.
\end{proof}

With a single gap, depth and exponent cutoffs are closely related. With
several gaps they need not agree. For instance, depth one for
$f=x+x^2+x^4$ gives $y-y^2-y^4$ while the omitted depth-two term $2y^3$ is larger
than the retained $y^4$ term. The depth-one approximation is still valid with
an $\OO(y^3)$ error, but it is not a complete exponent truncation through
$y^4$. This distinction is explicit in the returned metadata.

\subsection{Truncation by an absolute exponent of the output variable}

The public exponent cutoff $B$ means: retain every full block whose exponent
of $y$ is \emph{strictly less than} $B$. Since $t=(y/a)^{1/p}$, this is equivalent
to retaining
\[
 \Delta_k<D,\qquad D=pB-q.
\]
The package requires $B>q/p$, so the leading term is included.

Let
\[
 I_D=\{k\in\NN^m:\Delta_k<D\}.
\]
For $m>0$ its finite outer boundary is
\begin{equation}\label{eq:boundary}
 \partial I_D=\{k+e_j:k\in I_D,\ k+e_j\notin I_D,\ 1\le j\le m\}.
\end{equation}
Define
\begin{align}
 \Delta_*&=\min\{\Delta_k:k\in\partial I_D\},\label{eq:frontiermin}\\
 K_*&=\max\left\{\sum_jk_jK_j:
               k\in\partial I_D,\ \Delta_k=\Delta_*\right\}.
 \label{eq:frontierdegree}
\end{align}
These are computed from finitely many exact weights; no common denominator of
irrational exponents is introduced.

\begin{lemma}\label{lem:boundary}
$\Delta_*$ is the smallest weight of any omitted multiindex. Every omitted
multiindex of weight $\Delta_*$ is on $\partial I_D$.
\end{lemma}

\begin{proof}
Take a path from $0$ to an omitted index by adding coordinate unit vectors.
The first point on the path that leaves $I_D$ lies on the boundary, and all
subsequent steps strictly increase weight. Thus no omitted index can have
weight smaller than the boundary minimum. If an omitted index has this
minimum weight, subtracting any occupied coordinate must return to $I_D$;
otherwise there would be an omitted point of smaller weight. Hence that index
itself is a boundary point.
\end{proof}

\begin{theorem}[Exponent remainder]\label{thm:weighted}
For $G_q^{<B}(t)=\sum_{k\in I_D}t^{q+\Delta_k}Q_k^{(q)}(L)$,
\begin{equation}\label{eq:weightedbound}
 g(y)^q-G_q^{<B}(t)
 =\OO\left(t^{q+\Delta_*}T^{K_*}\right).
\end{equation}
The bound remains valid when different contributions at $\Delta_*$ cancel;
in that event it can overestimate the actual remainder.
\end{theorem}

\begin{proof}
Choose an integer $N$ with $(N+1)d>\Delta_*$. By \cref{thm:depth}, the infinite
tail after depth $N$ is
\[
 \OO\bigl(t^{q+(N+1)d}T^{(N+1)K}\bigr)
   =o(t^{q+\Delta_*}).
\]
There are finitely many remaining omitted indices with depth at most $N$.
Those of weight $\Delta_*$ have polynomial degree at most $K_*$ by
\cref{lem:boundary} and \eqref{eq:degreebound}. Every one with strictly larger
weight is $o(t^{q+\Delta_*})$ by \eqref{eq:power-beats-log}. This proves the
claimed bound. Cancellation only reduces the finite leading contribution.
\end{proof}

If the cutoff falls in a gap of the support, $\Delta_*>D$ and the actual
reported power is sharper than $B$. At a support point the logarithmic degree
must be retained: \eqref{eq:weightedbound} is not automatically
$\OO(y^B)$. The package's remainder is an \emph{asymptotic scale}, not an explicit
numerical constant or a computed interval of validity.

\section{The logarithmic example to arbitrary order}
\label{sec:logexample}

For $f(x)=x+x^2(1+\log x)$, take $a=p=1$, $m=1$, $\delta_1=1$, and
$P_1(L)=1+L$. Formula \eqref{eq:Qk} becomes
\begin{equation}\label{eq:logQ}
 Q_n(L)=\frac{(-1)^n}{n!}
       \prod_{r=1}^{n-1}(n+1+r+D_L)(1+L)^n,
 \qquad Q_0=1,
\end{equation}
and
\begin{equation}\label{eq:logseries}
 g(y)=\sum_{n\ge0}y^{n+1}Q_n(\log y).
\end{equation}
The first coefficient polynomials, generated independently by the verification
program, are
\begin{align}
Q_{0}(L) &= 1 \\
Q_{1}(L) &= - L - 1 \\
Q_{2}(L) &= 2 L^{2} + 5 L + 3 \\
Q_{3}(L) &= - 5 L^{3} - \frac{41 L^{2}}{2} - 27 L - \frac{23}{2} \\
Q_{4}(L) &= 14 L^{4} + \frac{241 L^{3}}{3} + \frac{335 L^{2}}{2} + 151 L + \frac{299}{6}

\end{align}
For example the third correction can also be checked directly:
\[
 -\frac16\frac{\dd^2}{\dd y^2}\bigl(y^6(1+\log y)^3\bigr)
 =-y^4\left(5(1+\log y)^3+\frac{11}{2}(1+\log y)^2
                                      +(1+\log y)\right).
\]
Expansion of this expression gives the fourth-power block in
\eqref{eq:log-answer}.

\subsection{Catalan numbers and harmonic numbers in the coefficients}

The highest logarithmic degree in $Q_n$ is exactly $n$. Its coefficient is
\begin{equation}\label{eq:Catalan}
 [L^n]Q_n(L)=\frac{(-1)^n}{n!}\prod_{s=n+2}^{2n}s
            =(-1)^n\frac{1}{n+1}\binom{2n}{n}
            =(-1)^n C_n.
\end{equation}
The empty product convention covers $n=1$. This identity supplies an
inexpensive high-order consistency check.

Let $H_j=\sum_{s=1}^j1/s$. The next coefficient is
\begin{equation}\label{eq:harmonic}
 [L^{n-1}]Q_n(L)=(-1)^n C_n\,n\bigl(1+H_{2n}-H_{n+1}\bigr),\qquad n\ge1.
\end{equation}
To prove it, the constant part of the operator product contributes the
coefficient $n$ of $L^{n-1}$ in $(1+L)^n$, while choosing one derivative
contributes an additional
$n\sum_{s=n+2}^{2n}1/s$. Two or more derivatives cannot contribute at this
degree. For $n=3$ the formula gives $-41/2$, agreeing with the explicit block.

More generally let $e_j$ denote an elementary symmetric polynomial, with
$e_0=1$. Then
\begin{equation}\label{eq:symmetric-Q}
 Q_n(L)=\frac{(-1)^n}{n!}
 \sum_{j=0}^{n-1}e_{n-1-j}(n+2,n+3,\ldots,2n)
       \frac{n!}{(n-j)!}(1+L)^{n-j}.
\end{equation}
This is a finite nonrecursive specification of every coefficient. It follows
by expanding the commuting operator product and differentiating $(1+L)^n$.

\subsection{The correct remainder after a displayed truncation}

For each fixed $N$,
\begin{equation}\label{eq:logremainder}
 g(y)=\sum_{n=0}^{N}y^{n+1}Q_n(\log y)
      +\OO\bigl(y^{N+2}|\log y|^{N+1}\bigr).
\end{equation}
In particular,
\[
 g(y)=y-y^2(1+\log y)+\OO(y^3|\log y|^2).
\]
Replacing this last remainder by $\OO(y^3)$ would be false, since the next
block has a nonzero $2\log^2y$ coefficient. This corrects the informal order
notation used in the original discussion without changing its leading two
terms.

\subsection{An explicit sufficient convergence condition and tail bound}

In this example the analytic argument can be made quantitative without solving
for a critical point. On the disk $|z-y|\le y/2$, $0<y<1$,
\[
 |z|\le\tfrac32y,\qquad
 |\log(z/y)|\le\sum_{j\ge1}\frac{(1/2)^j}{j}=\log2.
\]
Thus, with $A(y)=1+|\log y|+\log2$,
\[
 |h(z)|=|z^2(1+\log z)|\le\tfrac94 y^2A(y).
\]
Put $\rho(y)=9yA(y)$ and choose a homotopy radius
$\mathcal R=1/\rho(y)$. On $|\lambda|\le\mathcal R$,
$|\lambda h(z)|\le y/4<y/2$. The inverse root satisfies
$|g_\lambda(y)-y|<y/2$. Cauchy's estimate applied to $g_\lambda-y$ gives the
explicit bound
\begin{equation}\label{eq:explicit-log-tail}
 \boxed{\displaystyle
 \left|g(y)-\sum_{n=0}^{N}y^{n+1}Q_n(\log y)\right|
 \le\frac y2\frac{\rho(y)^{N+1}}{1-\rho(y)},\qquad \rho(y)<1.}
\end{equation}
This is sufficient, not optimal. It proves ordinary convergence of the
homotopy block sum for every positive $y$ satisfying the condition and yields
a computable error bound there. The sharper asymptotic order in
\eqref{eq:logremainder} follows immediately, although its best constant need
not resemble the conservative constant in \eqref{eq:explicit-log-tail}.

\section{The irrational-power example and its convergence radius}
\label{sec:powerexample}

Let $\alpha>1$, $d=\alpha-1$, and $f_\alpha(x)=x+x^\alpha$. Then
\begin{equation}\label{eq:alphaseries}
 g_\alpha(y)=\sum_{n\ge0}c_n y^{1+nd},\qquad
 c_0=1,
\end{equation}
where for $n\ge1$
\begin{align}
 c_n
 &=\frac{(-1)^n}{n!}\prod_{r=0}^{n-2}(n\alpha-r)\label{eq:alpha-product}\\
 &=\frac{(-1)^n\Gamma(n\alpha+1)}
         {\Gamma(n+1)\Gamma(nd+2)}
  =\frac{(-1)^n}{1+nd}\binom{n\alpha}{n}.
 \label{eq:alpha-gamma}
\end{align}
Formula \eqref{eq:alpha-product} follows directly from
$(-1)^nD_y^{n-1}y^{n\alpha}/n!$. It is preferable to a gamma quotient in symbolic
software because it is a finite polynomial in $\alpha$ for each $n$ and has no
apparent special-function singularities to simplify.

The first coefficients are
\begin{align*}
 c_1&=-1,&c_2&=\alpha,\\
 c_3&=-\frac{\alpha(3\alpha-1)}2,&
 c_4&=\frac{\alpha(4\alpha-1)(4\alpha-2)}6.
\end{align*}
For $\alpha=\sqrt2$ these give \eqref{eq:irrational-answer}; the next two are
\[
 c_5=-\frac{305}{12}+\frac{51\sqrt2}{4},\qquad
 c_6=-77+\frac{322\sqrt2}{5}.
\]
The powers are exactly $1+n(\sqrt2-1)$. Replacing $\sqrt2$ by a decimal number
would change the problem and would destroy exact identities between weights
in more complicated examples.

\subsection{A regular power series in the correct auxiliary variable}

Set $z=y^d$ and $u=g_\alpha(y)/y$. The equation is
\begin{equation}\label{eq:ualpha}
 u+zu^\alpha=1,
\end{equation}
so $u(z)=\sum_{n\ge0}c_nz^n$ is an ordinary analytic power series near $z=0$.
This is the right way to reduce the irrational-power problem to an analytic
implicit equation. Trying to force the original expansion into integer powers
of $y$ loses the natural scale.

\begin{proposition}\label{prop:radius}
The radius of convergence of $u(z)$ is
\begin{equation}\label{eq:radius}
 R_\alpha=\frac{(\alpha-1)^{\alpha-1}}{\alpha^\alpha}.
\end{equation}
The series \eqref{eq:alphaseries} therefore converges for
$0<y<R_\alpha^{1/(\alpha-1)}$. Its block coefficients satisfy
\begin{equation}\label{eq:stirling-coeff}
 |c_n|\sim\frac{\sqrt\alpha}{\sqrt{2\pi}(\alpha-1)^{3/2}}
 R_\alpha^{-n}n^{-3/2}.
\end{equation}
\end{proposition}

\begin{proof}
Apply Stirling's formula to the gamma quotient in
\eqref{eq:alpha-gamma}, with $d=\alpha-1>0$. The numerator is asymptotic to
$\sqrt{2\pi\alpha n}(\alpha n/\ee)^{\alpha n}$, while the denominator is
asymptotic to
$\sqrt{2\pi n}(n/\ee)^n\sqrt{2\pi}(dn)^{dn+3/2}\ee^{-dn}$.
Cancellation gives \eqref{eq:stirling-coeff}. The Cauchy--Hadamard root test
then gives \eqref{eq:radius} exactly.
\end{proof}

The associated critical point is also visible directly. Differentiating the
left side of \eqref{eq:ualpha} with respect to $u$ and setting it to zero gives
$1+\alpha zu^{\alpha-1}=0$. Together with the equation this yields
$u=\alpha/(\alpha-1)$ and $z=-R_\alpha$. This identifies the negative-real
singularity limiting the Taylor series in $z$, even though the positive real
inverse exists for all $y>0$.

For fixed $N$, the next exponent is $1+(N+1)d$, so the finite-order estimate is
particularly simple:
\[
 g_\alpha(y)=\sum_{n=0}^Nc_n y^{1+nd}
                +\OO\bigl(y^{1+(N+1)d}\bigr).
\]
There are no logarithmic factors to include in this example.

\section{Mixed scales, collisions, and nonunit leading powers}
\label{sec:examples}

\subsection{Two independent power corrections}

For
\[
 f(x)=x+A x^\alpha+B x^\beta,\qquad \alpha,\beta>1,
\]
put $d_1=\alpha-1$ and $d_2=\beta-1$. The depth-two expansion is
\begin{align}
 g(y)={}&y-Ay^\alpha-By^\beta
 +A^2\alpha y^{2\alpha-1}
 +AB(\alpha+\beta)y^{\alpha+\beta-1}
 +B^2\beta y^{2\beta-1}\notag\\
 &+\OO\bigl(y^{1+3\min(d_1,d_2)}\bigr).
 \label{eq:mixedpowers}
\end{align}
The error describes depth truncation, so some of the displayed higher-power
terms can be smaller than the error. For a complete ordered expansion one
instead chooses an exponent cutoff and includes every $k_1d_1+k_2d_2$ below it.

For $n=k_1+k_2\ge1$, the coefficient at the labelled multiindex is
\[
 \frac{(-1)^nA^{k_1}B^{k_2}}{k_1!k_2!}
 \prod_{r=0}^{n-2}(k_1\alpha+k_2\beta-r),
\]
multiplying $y^{1+k_1d_1+k_2d_2}$. This is the mixed generalized-binomial
analogue of \eqref{eq:alpha-product}. Polynomial logarithmic coefficients are
handled by replacing each scalar factor with the corresponding operator in
\eqref{eq:Qk}.

\subsection{An exact cancellation caused by a collision}

For $f=x+x^2+x^3$, the gaps are $1$ and $2$. The inverse starts
\begin{equation}\label{eq:collisionexample}
 g(y)=y-y^2+y^3-4y^5+14y^6-30y^7+33y^8+\OO(y^9).
\end{equation}
There is no $y^4$ term. At weight $3$, the multiindex $(3,0)$ contributes $-5$
and $(1,1)$ contributes $+5$. Both must be computed and added. A coefficient
engine that assigns one term to each input generator without cross products,
or that overwrites an earlier term at the same weight, fails this test.

The remainder estimator deliberately uses a support candidate before relying
on cancellation. It remains correct when the candidate coefficient vanishes;
it may then be conservative. Automatically seeking the first nonzero
coefficient beyond an arbitrary number of cancellations is not needed to
satisfy a finite requested cutoff.

\subsection{A ramified positive inverse}

For
\[
 f(x)=3x^2\bigl(1+x(1+\log x)\bigr),\qquad
 t=\sqrt{y/3},\quad P=1+\log t,
\]
the first terms are
\begin{equation}\label{eq:ramified}
 g(y)=t-\frac12t^2P
 +t^3\left(\frac58P^2+\frac14P\right)
 -\frac{t^4}{48}(6+D_L)(8+D_L)P^3
 +\OO(t^5T^4).
\end{equation}
The leading inverse is a square root, but the positive normalization makes the
branch unambiguous. The cutoff must still refer to the exponent of $y$, not of
$t$: retaining terms with $y$ exponent below $5/2$ retains the displayed
powers through $t^4$.

As an independent algebraic check, for $f(x)=3x^2(1+x^2)$ the exact positive
inverse is
\[
 g(y)=\sqrt{\frac{-1+\sqrt{1+4y/3}}{2}}
     =t\sqrt{\frac{2}{1+\sqrt{1+4t^2}}}.
\]
Expansion of the last ordinary analytic factor agrees with
\eqref{eq:Qk}. The verification program checks this identity through the
available finite order without using the same differential formula twice.

\subsection{Powers of the inverse}

For $f=x+x^2$, the inverse itself is $(\sqrt{1+4y}-1)/2$. Setting $q=2$ in
\eqref{eq:Qk} gives directly
\[
 g(y)^2=y^2-2y^3+5y^4-14y^5+\OO(y^6).
\]
This is a simple test of the factor $q/p^{|k|}$ in the general formula.
For arbitrary $p$, omitting that normalization factor would produce incorrect
coefficients even if all near-identity examples happened to pass.

\section{Residual certificates, forward remainders, and Newton refinement}
\label{sec:residual}

\subsection{From a forward residual to an inverse error}

\begin{proposition}[Residual transport]\label{prop:residual}
Suppose $g$ is the real inverse of $f$, $G$ is a positive approximation, and
$f'(x)\ge m>0$ on the interval between $G$ and $g(y)$. Then
\begin{equation}\label{eq:residual-bound}
 |G-g(y)|\le\frac{|f(G)-y|}{m}.
\end{equation}
\end{proposition}

\begin{proof}
Apply the mean value theorem to $f(G)-f(g(y))$.
\end{proof}

For the logarithmic example, the global constant $m_0$ in
\eqref{eq:global-slope} is valid for every positive approximation. For
$x+x^\alpha$, one may use $m=1$. These are rigorous analytic inequalities.
Evaluating either side with ordinary arbitrary-precision floating-point
arithmetic is still not the same thing as interval certification.

For the general model, an approximation with $G/t\to1$ and the true inverse
both lie in $[t/2,2t]$ eventually. From \cref{prop:localinverse}, eventually
\[
 f'(x)\ge ap\,2^{-1-|p-1|}t^{p-1}\qquad(t/2\le x\le2t).
\]
Thus
\begin{equation}\label{eq:scaledresidual}
 |G-g(y)|\le\frac{2^{1+|p-1|}}{ap}
 t^{1-p}|f(G)-y|
\end{equation}
for sufficiently small $t$. In particular, an inverse error of order
$t^{1+\sigma}T^K$ corresponds to a forward residual of order
$t^{p+\sigma}T^K$, not generally the same absolute exponent. This scaling is
important when $p\ne1$.

\subsection{Using a finite forward expansion of a more complicated function}

The package parses a finite expression exactly. A separate approximation to a
more complicated forward function may nevertheless be used with a justified
error budget.

\begin{proposition}[Transport of an omitted forward tail]\label{prop:forwardtail}
Let $f_0$ be an admissible finite model and suppose an actual real function
$\widetilde f$ satisfies, near zero,
\[
 \widetilde f(x)-f_0(x)=\OO(x^{p+\sigma}(1+|\log x|)^K),\qquad \sigma>0.
\]
Assume also that $\widetilde f'(x)\sim apx^{p-1}$ and that its positive inverse
$\widetilde g$ exists. Then
\begin{equation}\label{eq:forwardtail}
 \widetilde g(y)-g_0(y)=\OO(t^{1+\sigma}T^K).
\end{equation}
\end{proposition}

\begin{proof}
The hypotheses give $\widetilde g(y)\sim t$ and $g_0(y)\sim t$. Evaluate the
forward discrepancy at $g_0(y)$, obtaining a residual of order
$t^{p+\sigma}T^K$. Apply the scaled slope estimate to $\widetilde f$ on the
interval between its inverse and $g_0$.
\end{proof}

The derivative hypothesis must not be silently inferred by differentiating an
arbitrary Big-O estimate. Without derivative control, a tiny oscillatory
forward error can have a large derivative and can even spoil local
monotonicity. For an explicitly controlled analytic expansion such derivative
bounds may be available, but they are part of the input justification.

This proposition gives a principled way to preprocess, for example, an
analytic trigonometric factor by a sufficiently long Taylor polynomial. The
inverse's reported algorithmic remainder must then be combined with
\eqref{eq:forwardtail}. Increasing the inverse cutoff beyond the accuracy of
the supplied forward model does not create additional information about the
actual function.

\subsection{One Newton correction can improve an approximation}

Let $G\sim t$, set $r=f(G)-y$, and define
\[
 G_+=G-\frac{r}{f'(G)}.
\]
For admissible finite models, $f''(x)=\OO(t^{p-2})$ on $[t/2,2t]$ and
$f'(G)\asymp t^{p-1}$. Taylor's theorem therefore gives
\[
 G_+-g(y)=\OO\left(\frac{(G-g(y))^2}{t}\right).
\]
Hence an error $\OO(t^{1+\sigma}T^K)$ improves to
$\OO(t^{1+2\sigma}T^{2K})$, as long as the points remain in the local interval.
This is an analytic or numerical refinement; a symbolic Newton expression
must still be truncated in the correct scale to be called a canonical
power--log expansion.

\section{Algorithm, representation, and computational correctness}
\label{sec:algorithm}

\subsection{Normalization and the conservative parser}

The automatic parser first replaces the literal subexpression
\code{Log[x]} by a fresh symbol $L$ and expands the finite sum. Each summand
must be a product of explicit powers $x^\beta$ and factors independent of $x$;
the remaining coefficient must be a polynomial in $L$. Equal exponents that
can be proved equal are combined.

Next it identifies a provably least input exponent $p$. Its coefficient must
be a positive constant $a$, not a nonconstant polynomial in $L$. Every other
term is divided by $ax^p$, producing positive gaps $\delta_j$ and polynomials
$P_j$. Realness and positivity are checked under the supplied assumptions.
The parser does not have to decide an ordering between every pair of higher
symbolic input exponents merely to recognize a least leading term.

The grammar is intentionally narrower than mathematical equivalence under all
possible branch conditions. For example, the expression \code{Log[x\^{}2]}
is not automatically rewritten as \code{2 Log[x]}. A user who has already
proved the appropriate positive-domain identity can supply the normalized
expression or an explicit model. This keeps branch-dependent simplification
out of the parsing layer.

\subsection{Enumerating exactly the required multiindices}

Exponent-mode enumeration starts at $k=0$ and explores children $k+e_j$.
A child with weight below $D=pB-q$ is accepted and expanded in turn. A child at
or above the cutoff is recorded on the finite boundary but is not expanded.
A dictionary of integer multiindices avoids duplicate visits.

Every accepted index is reached because all its predecessors have smaller
weight. Every boundary candidate is reached from its accepted parent.
Conversely, no omitted subtree can re-enter the accepted region because all
gaps are positive. These statements prove the enumeration step independently
of any particular data structure.

For depth mode, the same traversal uses $|k|\le N$ instead of a weight
inequality. Its output is finite without ordering the symbolic gaps. The
number of accepted indices is $\binom{N+m}{m}$ before any normalizing
cancellations. Exponent-mode finiteness also follows from
$|k|<D/d$; a crude enclosing box has at most
$\prod_j(1+\lfloor D/\delta_j\rfloor)$ grid points.

\subsection{Exact weights and equality policy}

The implementation asks \code{FullSimplify} to prove equality, strict less
than, or strict greater than under the supplied assumptions. There is no
floating-point fallback. If an exponent-mode comparison cannot be resolved,
the function returns \code{Failure["UndecidableOrder",\ldots]} rather than
silently sorting approximate values.

In depth mode, a relative order may be unnecessary. Proven equalities are
collected, while unresolved possible coincidences remain an additive sum.
This representation is still algebraically valid; it does not claim the
weights are distinct. The metadata states that the ordering is by homotopy
depth rather than a canonical increasing-exponent list.

Machine and arbitrary-precision \code{Real} inputs are rejected, including
inexact coefficients as well as exponents. This rule is deliberately strict:
an exact identity between exponents is part of the mathematical input, not a
numerical tolerance decision. A user may choose an exact rational model
explicitly, but the package does not choose one on the user's behalf.

\subsection{Coefficient evaluation, collection, and reconstruction}

For each included index, compute $\Delta_k$, the polynomial product
$\prod P_j^{k_j}$, and the operator sequence in \eqref{eq:Qk}. Contributions
with provably equal weights are added before output. Proven zero polynomials
are removed. The reconstructed expression is
\begin{equation}\label{eq:reconstruct}
 \sum_k\left(\frac ya\right)^{(q+\Delta_k)/p}
 Q_k^{(q)}\!\left(\frac{\log(y/a)}p\right).
\end{equation}
This follows from the positive-domain normalization directly; the code does
not use \code{PowerExpand} to transform nested powers.

The result deliberately is not an ordinary \code{SeriesData} object.
\code{SeriesData} represents a fixed-denominator sequence of exponents, and
\code{InverseSeries} is documented as inversion of a power series
\cite{seriesdata,inverseseries}. Such a representation is not a universal
container for finite combinations of arbitrary irrational generators.
Moreover, a logarithmic error bound deserves explicit metadata even where
some built-in series operations can manipulate logarithmic coefficients.
The package does not redefine any built-in asymptotic function.

\subsection{Resource limits and the actual implementation cost}

The option \code{"MaxTerms"} bounds the number of visited labelled
multiindices, including boundary candidates. Exceeding it returns a failure;
no partial result is returned with an apparently complete remainder.

This is not a bound on all symbolic work. The current collector scans its
existing blocks for exact equality, so it can perform quadratically many
simplification queries in the number of blocks. Polynomial expression swell
and difficult parameter inequalities may also dominate execution time. The
source is intended as a transparent reference-quality implementation for
moderate finite orders, not as a claim of optimal high-order performance.
Canonical algebraic-number hashing, cached polynomial products, and a
sparse-vector implementation of \eqref{eq:vector-operator} are natural future
optimizations.

\begin{proposition}[Conditional computational specification]
If the implementation's exact real comparisons and polynomial operations
return their mathematically correct values, normalization accepts precisely
its stated grammar under the proved conditions, enumeration finishes within
its resource bound, and the output is constructed as described, then the
returned expression and remainder metadata satisfy \cref{thm:main,thm:depth,thm:weighted}.
\end{proposition}

\begin{proof}
Normalization gives \cref{def:model}. The traversal argument gives exactly the
required finite index set and its boundary. The coefficient routine implements
\eqref{eq:Qk}, and collection preserves the sum. Reconstruction is
\eqref{eq:reconstruct}. The cited remainder theorems apply to the recorded
truncation policy.
\end{proof}

This proposition is a mathematical specification and algorithmic correctness
argument, not a machine-checked proof of the Wolfram evaluator or an assertion
that native runtime testing has occurred.

\section{Wolfram Language package reference}
\label{sec:api}

The package context is \code{PowerLogInverse\textasciigrave}. Load the file from its
unpacked location; no installation, network request, or repository write is
performed by the package.
\begin{lstlisting}
Get[".../power-log-inverse/Kernel/PowerLogInverse.wl"];
Clear[x, y, ell];
\end{lstlisting}
Use distinct, unassigned symbols for the input variable, output variable, and
explicit logarithm variable.

\subsection{The main expression and data interfaces}

\begin{lstlisting}
PowerLogInverse[f, {x, y, B}, options]
PowerLogInverseData[f, {x, y, B}, options]
\end{lstlisting}
The first returns the explicit inverse approximation. The second returns an
association containing that expression, its normalization, branch conditions,
coefficient blocks, enumeration details, and remainder information. The
cutoff $B$ is a strict exponent of $y$, not a number of printed monomials.

The options are:
\begin{center}
\begin{tabularx}{\textwidth}{@{}p{34mm}p{25mm}X@{}}
\toprule
Option & Default & Meaning\\\midrule
\code{Assumptions} & \code{True} & Fixed-parameter conditions used to prove realness, positivity and exact comparisons.\\
\code{"Truncation"} & \code{"Exponent"} & Either strict output-exponent cutoff or \code{"Depth"}, which interprets the last argument as a nonnegative integer depth.\\
\code{"InversePower"} & $1$ & A fixed positive $q$; computes $g(y)^q$.\\
\code{"MaxTerms"} & $10000$ & Maximum visited multiindices, including the first rejected boundary.\\
\bottomrule
\end{tabularx}
\end{center}

The key fields of a successful result are:
\begin{longtable}{@{}p{46mm}p{101mm}@{}}
\toprule Field & Interpretation\\\midrule\endhead
\code{"Expression"} & The finite expression in the output variable.\\
\code{"NormalizedExpression"} & The same finite sum in fresh symbols $t,L$, before applying the positive normalization.\\
\code{"NormalizationRules"} & Rules $t\mapsto(y/a)^{1/p}$ and $L\mapsto\log(y/a)/p$.\\
\code{"Model"} & Validated leading coefficient, power, perturbations, logarithm variable, and assumptions.\\
\code{"Terms"} & List of weights, normalized exponents, exponents of $y$, and logarithmic polynomials.\\
\code{"Remainder"} & Either exact zero, or an association with \code{"Scale"}, \code{"YExponent"}, \code{"LogDegree"}, and the candidate omitted weight.\\
\code{"Conditions"} & Supplied model assumptions together with $y>0$.\\
\code{"Branch"} & \code{"PositiveNearZero"}.\\
\code{"IncludedMultiindices"} & Labelled indices used in the finite coefficient sum.\\
\code{"VisitedCount"} & Count used by the enumeration resource guard.\\
\code{"Ordering"} & Increasing exponent or a depth-mode qualification.\\
\bottomrule
\end{longtable}

A remainder field of type \code{"BigO"} means that the unknown error is bounded
by a constant times \code{"Scale"} for sufficiently small positive $y$, with
fixed admissible parameters. It does not identify the constant or the starting
point of that asymptotic regime. The main expression interface omits the
metadata for convenience; its mathematical domain remains the same.

\subsection{Explicit models and a single coefficient}

\begin{lstlisting}
m = PowerLogModel[a, p, {{delta1, P1}, {delta2, P2}}, ell,
      Assumptions -> conditions];
PowerLogInverse[m, {y, B}]
PowerLogInverseData[m, {y, B}]
PowerLogCoefficient[m, {k1, k2}]
\end{lstlisting}
Here $P_1,P_2$ are polynomials in the symbol \code{ell}. The model specifies
\eqref{eq:model}; it is not evaluated at a numerical $x$. The coefficient
function returns exactly \eqref{eq:Qk}, as a polynomial in the model's
logarithm variable. It also accepts \code{Assumptions} and
\code{"InversePower"}.

\begin{lstlisting}
m = PowerLogModel[1, 1, {{1, 1 + ell}}, ell];
PowerLogCoefficient[m, {6}]
\end{lstlisting}
This asks for the coefficient of $y^7$ in the logarithmic example. Automatic
normalization is also available separately as \code{PowerLogModel[f,x]}.
If repeated gaps combine or zero perturbations disappear, the length of the
multiindex must match the \emph{normalized} perturbation list stored in the
model, not an obsolete unnormalized input list.

\subsection{Descriptive failures}

Unsupported input returns a \code{Failure} rather than an arbitrary branch or
an unevaluated expression that looks like an answer. Important tags include
\code{"InexactInput"}, \code{"LogarithmicLeadingTerm"},
\code{"NonPolynomialLog"}, \code{"UnprovedPositive"},
\code{"UnprovedReal"}, \code{"UndecidableLeadingTerm"},
\code{"UndecidableOrder"}, \code{"CutoffTooLow"},
\code{"InvalidDepth"}, \code{"VariableCollision"}, and
\code{"TermLimit"}.

For instance,
\begin{lstlisting}
PowerLogInverse[x Log[x], {x, y, 3}]
PowerLogInverse[x + x^1.4, {x, y, 3}]
PowerLogInverse[x + x^2/Log[x], {x, y, 3}]
\end{lstlisting}
are intentionally rejected for different reasons. The first needs a
logarithmic leading core, the second supplies an inexact exponent, and the
third has a nonpolynomial logarithmic coefficient. No claim is made that the
underlying mathematical problems are impossible; they are outside this
package's stated finite polynomial-logarithm grammar.

\section{Worked package calls}
\label{sec:calls}

\subsection{The exact original use case}

\begin{lstlisting}
d = PowerLogInverseData[x + x^2 (1 + Log[x]), {x, y, 5}];
Collect[d["Expression"], y]
d["Remainder"]
\end{lstlisting}
The expression is \eqref{eq:log-answer} without its order term. The remainder
metadata has output exponent $5$ and logarithmic degree $4$, hence scale
$y^5(1+|\log y|)^4$.

For only the two leading blocks, replace $5$ by $3$. For arbitrary depth $N$,
use an exact integer $N$ and the depth option:
\begin{lstlisting}
PowerLogInverse[x + x^2 (1 + Log[x]), {x, y, 8},
  "Truncation" -> "Depth"]
\end{lstlisting}
This includes depths $0$ through $8$, or powers $y$ through $y^9$, not just
eight individual logarithmic monomials.

\Needspace{8\baselineskip}
\subsection{The irrational-power use case}

\begin{lstlisting}
PowerLogInverse[x + x^Sqrt[2], {x, y, 4 Sqrt[2] - 3}]
\end{lstlisting}
The cutoff is exactly the exponent of the first omitted term in the archive,
so the output contains its four requested terms. A generic $\alpha>1$ is
especially convenient in depth mode:
\begin{lstlisting}
PowerLogInverse[x + x^alpha, {x, y, 5},
  "Truncation" -> "Depth", Assumptions -> alpha > 1]
\end{lstlisting}
An exponent cutoff for symbolic $\alpha$ can require inequalities involving
integer multiples of $\alpha-1$; the appropriate number of included terms can
change with $\alpha$. The package does not silently choose one of those
parameter regions.

\subsection{Mixed exact irrational gaps}

\begin{lstlisting}
mixed = PowerLogInverseData[
  x + x^Sqrt[2] (1 + Log[x]) + 2 x^Sqrt[3],
  {x, y, 5/2}];
mixed["Terms"]
mixed["Remainder"]
\end{lstlisting}
Every weight generated by $\sqrt2-1$ and $\sqrt3-1$ below $3/2$ is included.
The logarithmic polynomial at each weight is produced by the same finite
operator formula. There is no conversion to a rational Puiseux denominator.

\subsection{Checking an approximation by a forward residual}

For the first example a numerical check can be performed as follows:
\begin{lstlisting}
value = 10^-6;
approx = N[d["Expression"] /. y -> value, 80];
residual = approx + approx^2 (1 + Log[approx]) - value;
N[Abs[residual]/(1 - 2 Exp[-5/2]), 30]
\end{lstlisting}
The expression before numerical evaluation is an upper bound on the true
absolute inverse error by \cref{prop:residual}. Its floating-point evaluation
is not an outward-rounded interval. A rigorous numerical certificate would
also enclose rounding errors in the evaluation of the residual and slope.

\section{Executed verification and reproducibility}
\label{sec:verification}

\subsection{What was actually checked}

The independent program \code{Verification/verify.py} ran successfully with
Python 3.13.5, SymPy 1.14.0, and mpmath 1.3.0. It records
\input{verification-count.tex}\ successful checks in
\code{verification-results.json} and \code{verification-report.txt}.

The exact checks include direct differentiation of the logarithmic homotopy
formula through depth ten; the falling-factorial formula for arbitrary
$\alpha$ through depth ten; the Catalan leading-log coefficient; formal
Picard inversion in independent marker variables; substitution into the
normalized forward equation; degree bounds; a ramified exact radical inverse;
coincident weights and their cancellation; and comparison of weighted
breadth-first enumeration with an independent bounded enumeration.

The formal-jet tests include nonunit leading power, rational and irrational
gaps, two logarithmic generators, and a positive noninteger power of the
inverse. Each jet operation uses ordinary finite polynomial multiplication,
truncated binomial powers, and the Taylor series of $\log(1+U)$. It does not
use the theorem being tested to construct the independent Picard solution.

\subsection{Numerical comparisons}

For each model, high-precision bisection on a positive bracket supplies a
reference inverse. The table uses 100 decimal digits and reports the actual
absolute error and the evaluated residual divided by the global slope lower
bound. The integer $N$ is the last retained perturbation depth; the sum includes
the leading depth-zero term. All entries are computed results, not fitted
error laws.

\begin{longtable}{@{}llrll@{}}
\caption{Independent numerical checks; floating point, not interval certification.}\label{tab:numerical}\\
\toprule Model & $y$ & $N$ & Absolute inverse error & Residual/slope bound\\\midrule\endfirsthead
\toprule Model & $y$ & $N$ & Absolute inverse error & Residual/slope bound\\\midrule\endhead
\tableinput{numerical-table.tex}
\bottomrule
\end{longtable}

The label \code{log} denotes \eqref{eq:original}, and \code{sqrt2} denotes
$x+x^{\sqrt2}$. For example, at $y=10^{-8}$ and depth five the errors are
approximately $3.07789\cdot10^{-47}$ and $1.84415\cdot10^{-27}$, respectively.
The much slower improvement in the second case is consistent with its small
gap $\sqrt2-1$ rather than a unit gap.

The JSON report also compares the signed error with the first omitted block.
For the logarithmic case at depth five and $y=10^{-8}$ that ratio is about
$1.0000005439$; for the irrational-power case it is about $0.9990508345$.
These observations support the leading-remainder analysis but do not replace
the proofs.

\subsection{What was not executed}

The connected Wolfram context/evaluator service returned HTTP 404, and no
local Wolfram kernel was available. Therefore the package's native runtime
behavior and the 35 supplied MUnit regression tests were \emph{not} executed
in this session. The source received a manual audit and a lexical
string/comment/delimiter check. That lexical checker is not a Wolfram parser
and is not evidence of runtime correctness.

The native regression suite covers the original examples, ordinary
\code{InverseSeries} cross-checks, explicit models, symbolic assumptions,
normalization, observables, collisions, remainder metadata, and expected
failures. Run it in a native installation before integrating the package into
a larger workflow:
\begin{lstlisting}[language={}]
wolframscript -file Tests/RunTests.wl
\end{lstlisting}
Alternatively evaluate \code{TestReport[".../Tests/PowerLogInverse.wlt"]} in a
notebook. The supplied command-line runner records the kernel version and
native report. It supports the current documented
\code{"ReportSucceeded"} property and checks available properties before
using the legacy success counts \cite{testreport,testreportobject,verificationtest}.
No fabricated native test report is included.

\subsection{Rebuilding the article and the verification outputs}

From the unpacked archive root, run
\begin{lstlisting}[language={}]
python Verification/verify.py
python Verification/static_check.py
cd Article
pdflatex -interaction=nonstopmode -halt-on-error power-log-inverse.tex
pdflatex -interaction=nonstopmode -halt-on-error power-log-inverse.tex
\end{lstlisting}
The verification program regenerates the numerical table and coefficient
blocks included by the article. The document uses standard TeX Live packages
and an inline bibliography; it needs no network access or bibliography
processor. The software license and article license are recorded separately
in the release.

\section{Limits of the implementation and principled extensions}
\label{sec:limits}

\subsection{A different leading core is a different problem}

A model such as $x\log(1/x)$ has a logarithmic leading factor, so it is not of
\cref{def:model}. Its positive inverse near zero uses the large negative
Lambert branch: writing $u=\log(1/x)$ gives $y=u\ee^{-u}$ and hence
\[
 x=-\frac{y}{W_{-1}(-y)}.
\]
One can then study perturbations around this exact logarithmic core. The
normalization is different from $t=(y/a)^{1/p}$ and requires a different
coefficient algebra. Rejecting the input in the present package is preferable
to disguising that change of problem as a power-series inversion.

Likewise, negative or fractional powers of $-\log x$, iterated logarithms,
and oscillatory coefficients can require a richer coefficient ring and new
error estimates. The polynomial logarithm condition is what makes every
coefficient here a finite, explicitly representable polynomial.

\subsection{The homotopy identity is more general than the parser}

The analytic identity \eqref{eq:homotopy} applies whenever its local analyticity
and smallness hypotheses hold. Some exponentially small perturbations or
other analytic germs on punctured domains satisfy those hypotheses even
though they are not finite power--log polynomials. The universal differential
formula can therefore be a starting point for further work.

What does not follow automatically is a useful ordering of the resulting
scales, a finite support below a requested cutoff, or the logarithmic degree
bound used by this implementation. Those are separate structural properties,
and this release does not advertise an unrestricted logarithmic-exponential
transseries engine.

\subsection{Fixed parameters versus uniform asymptotics}

All constants in the convergence and remainder proofs may depend on $p$, the
gaps, coefficients, and polynomial degrees. The estimates are not uniform as
$p\to0$, a gap tends to zero, or the number of generators grows. In particular,
when $\alpha\downarrow1$, the natural gap in the irrational-power example
shrinks and a fixed exponent cutoff contains more and more depths.

The proof explains the obstruction rather than hiding it: the local-finiteness
bound uses $d=\min\delta_j>0$, and the homotopy coefficient estimate uses
$t^dT^K\to0$ for fixed $d$. A simultaneous limit may require a new rescaled
leading equation.

\subsection{Useful next implementation steps}

The most valuable extensions are concrete: a native-kernel compatibility run
and continuous regression testing; cached sparse polynomial vectors;
canonical exact-algebraic weight keys; optional interval evaluation of
residual certificates; and a forward-expansion object carrying both function
and derivative remainder contracts as in \cref{prop:forwardtail}. Extending to
nonpolynomial logarithmic coefficients should come with a corresponding
remainder theorem rather than only an additional parser rewrite.

\section{Conclusion}

The requested positive inverse expansions are obtained by choosing the
correct endpoint germ and reverting a small perturbation of an explicit
monomial core. Logarithms and irrational powers then cease to be exceptional
cases: they are coefficients and weights in a locally finite algebra.

The central formula \eqref{eq:Qk} computes every requested coefficient by
finite polynomial operations. The homotopy argument proves that the formal
blocks represent the genuine positive inverse near zero, and the boundary
construction supplies valid logarithm-aware remainder metadata for exact
exponent cutoffs. The source package implements this specification without
silently rationalizing powers or choosing an arbitrary complex branch.

The archive includes the complete implementation, examples, native regression
tests, and the executed independent mathematical checks. The remaining
software verification limitation is explicit: the native Wolfram tests still
need to be run in an available kernel.

\appendix
\section{A minimal near-identity reference implementation}
\label{app:minimal}

For the two original examples, both of which have $f(x)=x+h(x)$ with $h=o(x)$,
the essential finite formula fits in a few lines. This educational routine
assumes the analytic hypotheses; it does not replace the validated parser,
weight collector, or remainder metadata of the main package.

\begin{lstlisting}
Clear[NearIdentityInverse];
NearIdentityInverse[f_, {x_Symbol, y_Symbol}, n_Integer] /;
    (n >= 0 && x =!= y) :=
 Module[{h = f - x},
  Expand[y + Total[Table[
    ((-1)^k/Factorial[k] D[h^k, {x, k - 1}]) /. x -> y,
    {k, 1, n}]]]
 ];

NearIdentityInverse[x + x^2 (1 + Log[x]), {x, y}, 3]
NearIdentityInverse[x + x^Sqrt[2], {x, y}, 3]
\end{lstlisting}
\code{Table} is used so that the differentiation order is an actual
nonnegative integer at each evaluation. This is a perturbation-depth formula.
For general mixtures it does not promise an exponent-complete cutoff, and it
provides no automatic branch or remainder checking. Its role is to expose the
simple mathematical core of the construction.

\section{Derivation of the polynomial operator identity}

For completeness, the general one-step rule is
\[
 \frac{\dd}{\dd s}\left[s^bZ\!\left(\frac{\log s}{p}\right)\right]
 =s^{b-1}\left(b+\frac1pD_L\right)Z(L).
\]
Apply it a second time to obtain
\[
 s^{b-2}\left(b-1+\frac1pD_L\right)
          \left(b+\frac1pD_L\right)Z(L).
\]
Induction yields \eqref{eq:derivative-identity}. The crucial point is that the
polynomial argument changes under differentiation; treating $\log s$ as a
constant would omit every $D_L$ term and produce wrong lower logarithmic
coefficients. For example it would retain the Catalan top coefficient in the
logarithmic example while giving $2L^2+4L+2$, rather than
$2L^2+5L+3$, in the cubic block. A test checking
only leading logarithms is therefore insufficient.

At fixed $k$, put $A_r=q+\Delta_k+pr$. Expanding the commuting product also
gives
\[
 Q_k^{(q)}(L)=\frac{(-1)^{|k|}q}{p^{|k|}k!}
 \sum_{j=0}^{|k|-1}e_{|k|-1-j}(A_1,\ldots,A_{|k|-1})
 D_L^j\!\left(\prod_iP_i(L)^{k_i}\right).
\]
This is an alternative finite coefficient formula. Both versions have the
same polynomial degree bound; the sequential-operator form generally avoids
constructing the elementary symmetric polynomials explicitly.

\section{Source provenance and scope of consultation}

The source of the precise question was the user-supplied archive, whose
Markdown text was read in full. The public question was also inspected to
confirm the branch-selection issue and the irrational-power example
\cite{question}. Historical command outputs in that question are treated as
historical observations, not as present-day benchmarks.

For the repository material, the directory listing, README, companion
introduction, selected coefficient/remainder discussion, synthesis, and Wolfram
recipe appendix were inspected through the GitHub connector. The large main
\emph{Transseries and Inversion} source was identified, but the available
content/blob endpoints did not return its text successfully. Consequently no
claim is made to have read that monograph in full. The present proofs are
self-contained and do not depend on inaccessible portions.

The official Wolfram documentation was consulted for the roles of
\code{InverseSeries}, \code{SeriesData}, \code{AsymptoticSolve}, and the test
framework \cite{inverseseries,seriesdata,asymptoticsolve,testreport,testreportobject,verificationtest}.
Native behavior of the archived commands was not rerun. In particular, no
blanket statement is made that current built-in asymptotic functions fail on
all expressions in the class studied here.

\clearpage
\begin{thebibliography}{99}
\bibitem{question}
Vladimir Reshetnikov,
\emph{How to get an asymptotic of the real-valued branch of the inverse function?},
Mathematica Stack Exchange, question 236367, December 11, 2020, with update
December 13, 2020. User-supplied archive and public page consulted September 7, 2026.
\url{https://mathematica.stackexchange.com/questions/236367/how-to-get-an-asymptotic-of-the-real-valued-branch-of-the-inverse-function}.

\bibitem{dlmf}
NIST Digital Library of Mathematical Functions,
\S1.10(iv), residue and Rouch\'e theorems, and \S1.10(vii), inverse functions
and Lagrange inversion. Consulted September 7, 2026.
\url{https://dlmf.nist.gov/1.10}.

\bibitem{repo}
Vladimir Reshetnikov / ProveIt repository,
\emph{Series and transseries}, directory README and file listing.
Consulted September 7, 2026.
\url{https://github.com/VladimirReshetnikov/ProveIt/tree/main/Analysis/FabiusFunction/docs/semi-formalized-research-frontiers/drafts/series-and-transseries}.

\bibitem{companion}
ProveIt repository,
\emph{Combinatorial coefficient algebras for transseries inverses}, research
draft, selected sections and appendices consulted September 7, 2026.
\url{https://github.com/VladimirReshetnikov/ProveIt/blob/main/Analysis/FabiusFunction/docs/semi-formalized-research-frontiers/drafts/series-and-transseries/Combinatorial_Transseries_Inverses/Combinatorial_Transseries_Inverses.tex}.

\bibitem{inverseseries}
Wolfram Research, \emph{InverseSeries}, Wolfram Language documentation.
Consulted September 7, 2026.
\url{https://reference.wolfram.com/language/ref/InverseSeries.html}.

\bibitem{seriesdata}
Wolfram Research, \emph{SeriesData}, Wolfram Language documentation.
Consulted September 7, 2026.
\url{https://reference.wolfram.com/language/ref/SeriesData.html}.

\bibitem{asymptoticsolve}
Wolfram Research, \emph{AsymptoticSolve}, Wolfram Language documentation.
Consulted September 7, 2026.
\url{https://reference.wolfram.com/language/ref/AsymptoticSolve.html}.

\bibitem{testreport}
Wolfram Research, \emph{TestReport}, Wolfram Language documentation.
Consulted September 7, 2026.
\url{https://reference.wolfram.com/language/ref/TestReport.html}.

\bibitem{testreportobject}
Wolfram Research, \emph{TestReportObject}, Wolfram Language documentation.
Consulted September 7, 2026.
\url{https://reference.wolfram.com/language/ref/TestReportObject.html}.

\bibitem{verificationtest}
Wolfram Research, \emph{VerificationTest}, Wolfram Language documentation.
Consulted September 7, 2026.
\url{https://reference.wolfram.com/language/ref/VerificationTest.html}.
\end{thebibliography}
\end{document}
